\textbf{Causal structure learning from IID data. }Causal structure learning has become an active area of research. hasan et al. \citep{hasan2023survey} recently presented a comprehensive review of causal discovery methods for IID data and time series. For IID data, several approaches rely on conditional independence to infer causal relationships from observational data, such as the classical PC algorithm \citep{spirtes2000causation}. Additionally, some methods extend beyond observational data, incorporating interventional data to enhance causal inference, including COmbINE \citep{triantafillou2015constraint} and HEJ \citep{hyttinen2014constraint}. These approaches utilize data collected from controlled interventions to uncover causal relationships. A novel research direction introduced in \citep{zheng2018dags} addresses the combinatorial challenges of structure learning by formulating it as a continuous constrained optimization problem, thus avoiding computationally costly combinatorial searches. Similarly, Zhu et al. \citep{zhu2019causal} utilize the acyclicity constraint but employ reinforcement learning techniques for estimating directed acyclic graphs (DAGs). In contrast, Ke et al. \citep{ke2019learning} propose an approach that learns DAGs from interventional data by optimizing an unconstrained objective function. \citep{brouillard2020differentiable} provide a comprehensive analysis of continuous-constrained methods, offering a generalized framework applicable to interventional data scenarios. Another significant method, DiBS \citep{lorch2021dibs}, estimates the full posterior distribution over Bayesian networks from limited observations, enabling quantification of uncertainty and assessment of confidence in causal discovery.

\textbf{Causal structure learning from stationary MTS. }The previously mentioned state-of-the-art methods primarily target independent observations rather than temporal dependencies. Assaad et al. \citep{assaad2022survey} provide a comprehensive review of approaches specifically designed for causal discovery from MTS. To model causal relationships involving time dependencies, researchers frequently employ Dynamic Bayesian Networks (DBNs), which effectively capture discrete-time temporal dynamics within directed graphical frameworks. Some methods neglect contemporaneous (instantaneous) dependencies and focus exclusively on recovering time-lagged causal links \citep{haufe2010sparse,song2009time}, and tsFCI \citep{entner2010causal}, the latter adapting the Fast Causal Inference algorithm \citep{spirtes2000causation} for time series data. Runge et al. \citep{runge2019detecting} introduced PCMCI, a scalable two-stage algorithm for time series, initially focusing only on time-lagged relationships. They subsequently extended it to PCMCI+ \citep{runge2020discovering}, enabling the identification of contemporaneous causal connections. Additionally, models addressing non-Gaussian instantaneous effects have been developed, such as VARLINGAM \citep{hyvarinen2010estimation}, which integrates non-Gaussian instantaneous models with autoregressive components. Another significant method is Time-series Models with Independent Noise (TiMINo) \citep{peters2013causal}, which studies nonlinear and instantaneous effects using constrained SEMs. Pamfil et al. \citep{pamfil2020dynotears} recently proposed DYNOTEARS, leveraging an algebraic characterization of graph acyclicity from \citep{zheng2018dags} to estimate both instantaneous and time-lagged relationships from time series data. DYNOTEARS utilizes a score-based DBN learning approach optimized via an augmented Lagrangian framework, enabling causal graph inference without assumptions on the underlying topology. In contrast, methods like NBCB \citep{assaad2021mixed}, a noise-based/constraint-based approach, aim to learn a summary causal graph directly from observational time series data, going beyond Markov equivalence constraints even in the presence of instantaneous relationships. Rhino \citep{gong2022rhino} introduced the first model that tackle stationary MTS with historical noise,where they assume that the noise variance changes over time as a function of solely time lagged parents. All the aforementioned methods does not handle heteroscedastic setting and also fail short in the case of MTS with multiple regimes.

\textbf{Causal structure learning from MTS with multiple regimes. }Some research have aimed to address this challenge by developing methods for causal discovery in heterogeneous data. An example of such a method is CD-NOD \citep{huang2020causal}, tackles time series with various regimes. By using the time stamp IDs as a surrogate variable, CD-NOD output one summary causal graph where the parents of each variable are identified as the union of all its parents in graphs from different regimes. Then it detects the change points by using a non-stationary driving force that estimates the variability of the conditional distribution $p(x_i|\text{union parents of } x_i)$ over the time index surrogate. While CD-NOD provides a summary graph capturing behavioral changes across regimes, it falls short in inferring individual causal graphs, also CD-NOD cannot handle either non-Gaussian or heteroscedastic noise. The overall summary graph does not effectively highlight changes between regimes. Additionally, CD-NOD detects the change points but fails to determine the regime indices, rendering it incapable of inferring the precise number of regimes. In scenarios involving recurring regimes, CD-NOD is unable to detect this crucial information. Another relevant work dealing with MTS composed of multiple regimes is RPCMCI \citep{saggioro2020reconstructing}. In this paper, the model \citep{saggioro2020reconstructing} learns a temporal graph for each regime. However, they focus initially on inferring only time-lagged relationships and require prior knowledge of the number of regimes and transitions between them. \citep{balsellsidentifiability} addresses first-order regime-dependent causal discovery from MTS with multiple regimes. They proved that first-order Markov switching models with non-linear Gaussian transitions are identifiable up to permutations. Their work offers also a practical algorithms for regime-dependent causal discovery in time series data. However, its primary limitation is the assumption of solely time-lagged relationships, with the theory being restricted to a single time lag. CASTOR \citep{rahmanicausal} learns regime indices and their corresponding causal graphs, including instantaneous and time-lagged relationships, under the assumption of normally distributed noise with equivariance. However, like other causal discovery methods for non-stationary MTS, it does not address non-Gaussian or heteroscedastic noise. \citep{varambally2023discovering} tackle a different setting in which they aim to discover a mixtures of Structural Causal Models from a datasets of MTS. They assume that they have different stationary MTS in the same dataset, one regime per MTS, and each one could be explained by one causal model in the mixture. In their case, they assume that every MTS in the dataset is stationary but the whole dataset is a mixture. In our case, we assume that we have only one non-stationary MTS and it is composed of different regimes where we do not know when the regime starts and ends, and our goal is to identify the regimes and the corresponding causal graphs.